% ----------------------------------
%   - Derivation of the Black-Scholes PDE
%       - Assumptions
%       ? Rigorous construction of market bla bla
%       - From SDE to PDE
%   - Solving the Black-Scholes PDE (BS Formula)
% ----------------------------------

\section{Black-Scholes}

Black-Scholes is a mathematical model for pricing options.
Although many other models have been developed since,
the Black-Scholes model remains as the most important model in option pricing,
as in many cases options in markets are quoted using the Black-Scholes formula.

A derivative is a financial instrument whose payments are derived from the price of an other asset, called the underlying asset.
Derivatives are widely used in finance,
whether to hedge financial risk, this is to protect against adverse price movements,
or to bet on specific views, especially when it's difficult to take a position in the underlying asset.
So being able to price them and justify the price is key for the market to function.

\subsection{Replication and arbitrage}

One of the most important concepts in modern finance is the so called \textit{law of one price}.
The idea is that if two financial assets have exactly the same payoffs in all possible scenarios,
then they should have the same price.

This is a linked to the concept of \textit{arbitrage},
informally we say that arbitrage is the opportunity to make a riskless profit.
For example, if two assets that have exactly the same payoffs in all possible scenarios,
but don't follow the law of one price,
one can buy the cheaper one and sell the more expensive one,
gaining and instant profit and compensating the payoffs of the sold asset using the payoffs of bought one,
which match exactly.

We will work under the assumption that we are in operating in a market with no arbitrage,
this assumption is technically known as \textit{The First Fundamental Theorem of Asset Pricing},
and a key part of the \textit{Efficient Market Hypothesis}.
Under this assumption, we can price any derivative by creating a replicating portfolio,
this is a set of assets that has exactly the same payoffs as the derivative in all possible scenarios.

There are two main ways of creating a replicating portfolio:
\begin{itemize}
    \item \textbf{Static replication} refers to an replicating portfolio that is created at the beginning and held as is until maturity. 
    \item \textbf{Dynamic replication} refers to a replicating portfolio that requires a continuous change of its composition over time.
\end{itemize}

For a further discussion on replication, and its applications in finance, we recommend the book \cite{derman_vol_smile}.

% Weak and strong replication?

Will show two examples of static replication from \cite{steele_stochastic}
that will be useful to understand the Black-Scholes model.

\subsubsection{Example: Forward contract}

Assume that there is an agent that borrows and lends money at a continuously compounded rate $r$, known as the \textbf{risk-free rate}.
This means that if we borrow $1 \$$ at time $t = 0$,
we will have to pay $e^{rT} \$$ at time $t = T$.

A forward contract gives the holder the obligation to buy an asset at a fixed price $K$ at a fixed date $T$.

To price this contract imagine we are at time $t$ and the current price of an asset is $S$,
we want to buy a forward contract with maturity $T$ and fixed price $K$.

We have two options, we can buy the contract at time $t$ for a price of $F$ and at time $T$ pay $K$ to get the asset.
Alternatively we can buy the asset for $S$ and borrow $e^{-r(T - t)}K$ at time $t$, so at time $T$ we will pay $K = e^{r(T - t)} e^{-r(T - t)} K$ while having the asset.
Both options have the same cashflow at time $T$, pay $K$ and own the asset,
so either we have
\begin{equation*}
    F = S - e^{-r(T - t)} K,
\end{equation*}
or there is an arbitrage opportunity.

We have found the price of the forward contract by replicating its cashflows, using the underlying assed and the risk-free rate,
plus arguing that there should not be an arbitrage opportunity.

\subsubsection{Example: Put-Call parity}

\begin{definition}[European call option]
    An \textbf{European call option} is a financial contract that gives the holder the right,
    but not the obligation,
    to \textbf{buy} an asset at a fixed price (the strike price) at a fixed date (the maturity date).

    The payoff of an European call option at maturity $T$ is given by:
    \begin{equation}
        C_T = \max(S_T - K, 0),
    \end{equation}
    where $S_T$ is the price of the underlying asset at maturity and $K$ is the strike price.

    \begin{figure}[h] \label{fig:call_payoff}
        \centering
        \begin{tikzpicture}[scale=1.0, thick]
            % Axis
            \draw[->] (0,0) -- (7,0) node[anchor=north] {$S_T$};
            \draw[->] (0,-1) -- (0,4) node[anchor=east] {$C_T$};
        
            % Payoff line
            \draw[blue, thick, domain=0:3] plot(\x, {0});
            \draw[blue, thick, domain=3:6] plot(\x, {\x - 3});
        
            % Labels
            \node at (3,-0.3) {$K$};
            \node at (-0.3,0) {0};
            % \node at (-0.8,2) {$S_T - K$};
        
        \end{tikzpicture}
        \caption{Payoff of a European call option}
    \end{figure}
\end{definition}

\begin{definition}[European put option]
    An \textbf{European put option} is a financial contract that gives the holder the right,
    but not the obligation,
    to \textbf{sell} an asset at a fixed price (the strike price) at a fixed date (the maturity date).

    The payoff of an European put option at maturity $T$ is given by:
    \begin{equation}
        P_T = \max(K - S_T, 0),
    \end{equation}
    where $S_T$ is the price of the underlying asset at maturity and $K$ is the strike price.

    \begin{figure}[h] \label{fig:put_payoff}
        \centering
        \begin{tikzpicture}[scale=1.0, thick]
            % Axis
            \draw[->] (0,0) -- (7,0) node[anchor=north] {$S_T$};
            \draw[->] (0,-1) -- (0,4) node[anchor=east] {$C_T$};
        
            % Payoff line
            \draw[blue, thick, domain=0:3] plot(\x, {3 - \x});
            \draw[blue, thick, domain=3:6] plot(\x, {0});
        
            % Labels
            \node at (3,-0.3) {$K$};
            \node at (-0.3,0) {0};
            % \node at (-0.8,2) {$S_T - K$};
        
        \end{tikzpicture}
        \caption{Payoff of a European put option}
    \end{figure}
\end{definition}

Now consider a portfolio where we buy a call option and sell a put option with same strike $K$ and maturity $T$,
the payoff of such portfolio at time $T$ will be exactly $S_T - K$.
So we are essentilly reciving the asset and paying $K$ at time $T$,
as we have seen this is the payoff of a forward contract, which means that at time $t$ the price of our portfolio must be equal to the price of a forward of the asset,
\begin{equation} \label{eq:put_call_parity}
    C_t - P_t = F_t = S_t - e^{-r(T - t)} K,
\end{equation}
this formula is called \textit{put-call parity} and illustrates the relationship between the price of a call and a put option.


\subsection{The Black-Scholes model}\label{subsec:black_scholes_model}

The Black-Scholes model is a mathematical model for pricing options,
that is based on replicating the payoffs of the option dynamically using the underlying asset and a risk-free rate bond.

Like before we assume that the market is arbitrage-free,
and that there is a constant risk-free rate $r$.
As the model is based on dynamic replication,
we also need to assume that we can rebalance our replicating portfolio continuously,
i.e. we can buy and sell any amount of the underlying asset and the risk-free rate bond instantaneously at any time.

Let $S_t$ be the price of the underlying asset at time $t$ and $\beta_t$ be the price of the risk-free rate bond at time $t$.
The dynamics of this two assets are given by the following stochastic differential equations (SDEs):
\begin{equation*}
    dS_t = \mu S_t dt + \sigma S_t dW_t,
    \qquad
    d\beta_t = r \beta_t dt.
\end{equation*}

To construct the replicating portfolio.
Let $a_t$ be the amount of the underlying asset in our replicating portfolio and $b_t$ be the amount of the risk-free rate bond in our replicating portfolio.
Then the value of the portfolio at any time $t$ is given by:
\begin{equation*}
    V_t = a_t S_t + b_t \beta_t.
\end{equation*}

We need this portfolio to have the same payoff as the option at time $T$, so
\begin{equation*}
    V_T = \max{S_T - K, 0}.
\end{equation*}

As the portfolio is dynamically replicating the option,
the cashflow of the portfolio must match the ones of the option,
these are the payoff as seen before, and the initial payment,
so no cash must came in or out of the portfolio at any time between $t$ and $T$.
This means that the portfolio must be \textit{self-financing},
so the rebalanced must come from selling the underlying asset and buying the bond, or vice versa.
This is expressed mathematically as
\begin{equation}\label{eq:self_financing}
    dV_t = a_t dS_t + b_t d\beta_t,
\end{equation}
which can be extended to
\begin{equation*}
    dV_t = \left( a_t \mu S_t + b_t r \beta_t \right) dt + a_t \sigma S_t dW_t.
\end{equation*}

Now as the value of the replicanting portfolio $V_t$ should be equal to the value of the option at any time $t$,
we assume that this can be expressed by a function $f(t, S_t)$,
depending only on the time $t$ and the price of the underlying asset $S_t$.
Thus applying Ito's formula we have that
\begin{equation*}
    dV_t =
        \left( \frac{\partial f}{\partial t} + \frac{\partial f}{\partial S} \mu S_t + \frac{1}{2} \frac{\partial^2 f}{\partial S^2} \sigma^2 S_t^2 \right) dt
        + \frac{\partial f}{\partial S} \sigma S_t dW_t,
\end{equation*}
matching the two expressions for $dV_t$ we have that
\begin{equation*}
    a_t = \frac{\partial f}{\partial S},
\end{equation*}
and 
\begin{equation*}
    b_t =
        \frac{1}{r \beta_t}
        \left( \frac{\partial f}{\partial t} + \frac{1}{2} \frac{\partial^2 f}{\partial S^2} \sigma^2 S_t^2 \right)
\end{equation*}
as $V_t$ is the value of the replicating porfolio and the option at time $t$,
we have that
\begin{equation*}
    f(t, S_t) = V_t = a_t S_t + b_t \beta_t,
\end{equation*}
using the values of $a_t$ and $b_t$ from above,
\begin{equation*}
    f(t, S_t) =
        \frac{\partial f}{\partial S} S_t 
        + \frac{1}{r} \left( \frac{\partial f}{\partial t} + \frac{1}{2} \frac{\partial^2 f}{\partial S^2} \sigma^2 S_t^2 \right),
\end{equation*}
that can be rearranged to
\begin{equation}\label{eq:bs_pde}
    \frac{\partial f}{\partial t} =
        r f(t, S_t)
        - r \frac{\partial f}{\partial S} S_t 
        - \frac{1}{2} \frac{\partial^2 f}{\partial S^2} \sigma^2 S_t^2,
\end{equation}
with terminal condition
\begin{equation*}
    V_T = f(T, S_T) = \max(S_T - K, 0),
\end{equation*}
which is the \textbf{Black-Scholes PDE}.

